\documentclass[11pt,a4paper]{article}

% Paket untuk margin ATS-friendly
\usepackage[left=0.75in, right=0.75in, top=0.75in, bottom=0.75in]{geometry}
\usepackage{enumitem}
\usepackage{hyperref}
\usepackage{titlesec}

% Pengaturan hyperref agar link tidak berkotak
\hypersetup{
    colorlinks=true,
    linkcolor=black,
    filecolor=black,      
    urlcolor=black,
    pdftitle={CV Mulia Andiki},
}

% Pengaturan spasi antar paragraf dan penghapusan indentasi
\setlength{\parindent}{0pt}
\setlength{\parskip}{4pt}

% Pengaturan format judul bagian (Section)
\titleformat{\section}{\large\bfseries\uppercase}{}{0em}{}[\titlerule]
\titlespacing{\section}{0pt}{12pt}{8pt}

\begin{document}

% ================= HEADER =================
\begin{center}
    {\huge \textbf{MULIA ANDIKI}} \\ \vspace{4pt}
    muliaandiki@gmai.com | +62 82216903771 | \href{https://www.linkedin.com/in/mulia-andiki-030457331}{LinkedIn} | \href{https://github.com/MuliaAndiki}{GitHub} | \href{https://fortopolio-nu.vercel.app}{Portfolio}
\end{center}

% ================= RINGKASAN PROFESIONAL =================
\section*{Ringkasan Profesional}
Mahasiswa tingkat akhir di program studi Informatika yang berdedikasi sebagai Software Engineer dengan rekam jejak kuat pada Frontend Development serta minat mendalam terhadap Backend Development dan Full Stack Development. Memiliki pengalaman ekstensif dalam mengarsiteki solusi berbasis kecerdasan buatan (Artificial Intelligence) dan Internet of Things (IoT) untuk menyelesaikan masalah dunia nyata secara efisien. Terbiasa merancang platform modern berkinerja tinggi, mengimplementasikan infrastruktur Cloud Computing, dan memecahkan tantangan teknis pada Sistem Terdistribusi. Berorientasi pada Research and Innovation dengan komitmen untuk terus mendorong batasan teknologi melalui kolaborasi tim dan eksekusi proyek berbasis riset.

% ================= PENDIDIKAN =================
\section*{Pendidikan}
\textbf{Universitas Syah Kuala} \hfill Banda Aceh, Indonesia \\
S1 Informatika (Fakultas Matematika dan Ilmu Pengetahuan Alam) \hfill 2023 -- Sekarang (Semester 6) 

% ================= KEAHLIAN TEKNIS =================
\section*{Keahlian Teknis}
\begin{itemize}[leftmargin=0.15in, label={}]
    \item \textbf{Programming Languages:} JavaScript, TypeScript, Python.
    \item \textbf{Backend:} Bun, ElysiaJS, Express.js, FastAPI, Node.js.
    \item \textbf{Frontend:} Next.js, React, Tailwind CSS.
    \item \textbf{Database:} PostgreSQL, Monggodb.
    \item \textbf{DevOps \& Cloud:} Docker, Git, GitHub , GitAction
    \item \textbf{AI \& Machine Learning:} YOLOv8, Convolutional Neural Networks (CNN), TF-IDF 
    \item \textbf{IoT:} Arduino R4 Wifi
\end{itemize}

% ================= PENGALAMAN PROYEK =================
\section*{Pengalaman Proyek}

\textbf{Etno} \hfill Jan 2026 - Sekarang \\
\textit{Full Stack Developer} \\
Platform Pembelajaran Bahasa Indonesia Berbasis 3D Module 
\begin{itemize}[leftmargin=0.25in]
    \item \textbf{Tanggung Jawab:} Mengelola siklus pengembangan aplikasi dari sisi frontend hingga arsitektur backend dan intergrasi devops
    \item \textbf{Aksi:} Mengembangkan arsitektur aplikasi menggunakan Next.js dan Tailwind CSS pada sisi klien, serta membangun API server berkinerja tinggi menggunakan Bun dan ElysiaJS dengan PostgreSQL sebagai sistem manajemen basis data untuk menghasilkan sistem yang skalabel dan responsif.
\end{itemize}

\textbf{GETSMART} \hfill Mei 2026 - Sekarang \\
\textit{Software Engineer} \\
Platform Pembelajaran Berbasis Computer Vison Untuk Mendeteksi Emosi Pada Interface Muka
\begin{itemize}[leftmargin=0.25in]
    \item \textbf{Tanggung Jawab:} Frontend Developer Slicing Uiux Serta Mengintergrasikan API.
    \item \textbf{Aksi:} Mengintegrasikan layanan API untuk dasbor siswa, guru, dan admin menggunakan Google OAuth untuk memfasilitasi akses terpusat yang aman bagi berbagai peran pengguna.
\end{itemize}

\textbf{AERIS} \hfill April 2026 -- Mei 2026 \\
\textit{Software Engineer} \\
Dasbor lingkungan untuk memvisualisasikan kualitas udara, cuaca, dan risiko bencana secara real-time.
\begin{itemize}[leftmargin=0.25in]
    \item \textbf{Tanggung Jawab:} Merancang skema basis data dan menghubungkan data pihak ketiga ke dalam sistem.
    \item \textbf{Aksi:} Mengimplementasikan integrasi API eksternal menggunakan Open-Meteo dan API-Ninjas untuk menyajikan visualisasi data cuaca dan metrik lingkungan yang real-time.
\end{itemize}

\textbf{NutriPlate} \hfill Mei 2025 -- April 2026 \\
\textit{Software Engineer (IoT \& PWA)} \\
Sistem Nutritional Calculation berupa platform IoT/PWA untuk klasifikasi gizi anak menggunakan model visi komputer.
\begin{itemize}[leftmargin=0.25in]
    \item \textbf{Tanggung Jawab:} Mengintegrasikan lapisan frontend dan backend serta manajemen model AI.
    \item \textbf{Aksi:} Menggabungkan model Computer Vision ke dalam alur kerja aplikasi menggunakan YOLOv8 untuk menghasilkan kemampuan pemrosesan klasifikasi makanan anak yang responsif dan akurat.
\end{itemize}

\textbf{Loka - Loka} \hfill Jun 2025 - Sekarang \\
\textit{Software Engineer} \\
Sistem transaksi digital yang mendukung gateway pembayaran berlapis.
\begin{itemize}[leftmargin=0.25in]
    \item \textbf{Tanggung Jawab:} Mengonfigurasi dan memverifikasi integrasi pembayaran komersial.
    \item \textbf{Aksi:} Membangun saluran pemrosesan pembayaran menggunakan Midtrans dan GoPay QRIS untuk memastikan penyelesaian transaksi yang cepat dan validasi dokumentasi bisnis yang terjamin.
\end{itemize}

\textbf{KostHub} \hfill Maret 2025 - Sekarang \\
\textit{Software Engineer} \\
Platform manajemen pencarian dan penyewaan kos yang memudahkan interaksi antara pemilik properti dan penyewa melalui dasbor yang intuitif.
\begin{itemize}[leftmargin=0.25in]
    \item \textbf{Tanggung Jawab:} Bertanggung jawab atas pengembangan fitur pencarian, sistem manajemen properti, dan integrasi backend untuk alur data penyewaan.
    \item \textbf{Aksi:} Mengembangkan antarmuka pengguna yang responsif menggunakan Next.js dan membangun API RESTful yang aman menggunakan Express.js, serta melakukan integrasi database untuk menghasilkan sistem pengelolaan data kos yang efisien dan mempermudah proses transaksi penyewaan.
\end{itemize}

\textbf{CRM System} \hfill Maret 2025 - Agust 2025 \\
\textit{Software Engineer} \\
Sistem Manajemen Hubungan Pelanggan (CRM) yang dirancang untuk mengotomatisasi alur kerja interaksi klien dan memusatkan data prospek bisnis.
\begin{itemize}[leftmargin=0.25in]
    \item \textbf{Tanggung Jawab:} Mengembangkan modul fungsional sistem, mengoptimalkan kueri basis data untuk pemrosesan data prospek, dan memastikan integrasi layanan pihak ketiga yang aman.
    \item \textbf{Aksi:} Mengembangkan sistem pelacakan interaksi pelanggan menggunakan arsitektur backend yang tangguh dan teknologi [Sebutkan Teknologi Backend/Frontend, misal: Node.js/React], serta menerapkan protokol keamanan untuk menghasilkan sistem CRM yang mampu meningkatkan retensi pelanggan dan efisiensi operasional tim penjualan.
\end{itemize}


% ================= PENGALAMAN PENELITIAN DAN INOVASI =================
\section*{Pengalaman Dan Inovasi}

\textbf{NutriPlate -- Kompetisi Innovillage} \hfill 2025 -- 2026 \\
\textit{Software Engineer (IoT \& PWA)}
\begin{itemize}[leftmargin=0.25in]
    \item \textbf{Pencapaian:} Berhasil meraih posisi \textbf{Top 180} tingkat nasional dalam kompetisi inovasi sosial digital, Innovillage.
    \item \textbf{Deskripsi:} Platform IoT/PWA terintegrasi kecerdasan buatan (Sistem Nutritional Calculation) yang memanfaatkan visi komputer untuk klasifikasi gizi anak.
    \item \textbf{Kontribusi Utama:} Mengintegrasikan model machine learning (Computer Vision) ke dalam alur kerja aplikasi web progresif dan mengelola arsitektur backend untuk memastikan pemrosesan klasifikasi makanan berjalan dengan akurasi dan performa tinggi.
    \item \textbf{Teknologi yang digunakan:} YOLOv8, PWA (Progressive Web App), Integrasi Perangkat IoT, ElysiaJs, Fast Api.
\end{itemize}

% ================= ORGANISASI DAN KEPEMIMPINAN =================
\section*{Organisasi dan Kepemimpinan}

\textbf{Badan Eksekutif Mahasiswa (BEM) FMIPA, Universitas Syah Kuala} \hfill Banda Aceh, Indonesia \\
\textit{Anggota Departemen Publikasi, Dekorasi, dan Dokumentasi (PDD) - Videografer} \hfill 2024
\begin{itemize}[leftmargin=0.25in]
    \item Bertanggung jawab dalam perencanaan, pengambilan gambar, dan penyuntingan video untuk kebutuhan publikasi kegiatan serta pendokumentasian acara fakultas.
    \item Mengelola alur produksi konten visual kreatif guna meningkatkan jangkauan dan citra organisasi di media sosial.
\end{itemize}

\textbf{Himpunan Mahasiswa Informatika (Hmif), Universitas Syah Kuala} \hfill Banda Aceh, Indonesia \\
\textit{Anggota Himpunan} \hfill 2025 -- Sekarang
\begin{itemize}[leftmargin=0.25in]
    \item Berpartisipasi aktif dalam dua periode kepengurusan untuk mendukung program kerja himpunan serta memperkuat kolaborasi antar mahasiswa di lingkungan departemen.
\end{itemize}

\textbf{Penelitian dan Pengembangan Mahasiswa (PPM), Universitas Syah Kuala} \hfill Banda Aceh, Indonesia \\
\textit{Anggota} \hfill 2025 -- 2026
\begin{itemize}[leftmargin=0.25in]
    \item Berkontribusi sebagai anggota selama dua periode kepengurusan (kabinet yang berbeda), berfokus pada fungsi koordinasi, analisis, dan advokasi untuk meningkatkan kualitas serta inovasi di bidang kemahasiswaan.
\end{itemize}

% ================= PRESTASI DAN PENGHARGAAN =================
\section*{Prestasi dan Penghargaan}
\begin{itemize}[leftmargin=0.25in]
    \item \textbf{Top 180 Nasional, Innovillage} -- Berhasil menembus peringkat Top 180 tingkat nasional pada kompetisi inovasi sosial digital melalui pengembangan proyek NutriPlate (Sistem Nutritional Calculation berbasis IoT/PWA).
\end{itemize}

% ================= BAHASA =================
\section*{Bahasa}
\begin{itemize}[leftmargin=0.25in]
    \item \textbf{Bahasa Indonesia:} Native/Bilingual Proficiency
    \item \textbf{Bahasa Inggris:} Basic Proficiency
\end{itemize}

\end{document}